\documentclass[11pt,a4paper]{article}
\usepackage[utf8]{inputenc}
\usepackage[T1]{fontenc}
\usepackage[spanish]{babel}
\usepackage{amsmath}
\usepackage{amssymb}
\usepackage{geometry}
\geometry{margin=2.5cm}

\title{\textbf{Introducción a la Cosmología \\ Guía 5: Recombinación}}
\date{}

\begin{document}
\maketitle

\section*{Problema 1}
Asumiendo que la recombinación del hidrógeno es un proceso de equilibrio, estime la temperatura y el redshift a los que se produce. Observe que la temperatura de recombinación es muy inferior a la energía de ligadura del hidrógeno. ¿Por qué?

\section*{Problema 2}
Las observaciones cosmológicas indican que el universo es espacialmente plano, que su ritmo actual de expansión es $H_0\simeq 70\,(\text{km/s})/\text{Mpc}$ y que actualmente está compuesto en un 70\% de energía oscura (constante cosmológica) y en un 30\% de materia no relativista, de la cual una sexta parte es materia ordinaria. Teniendo esto en cuenta, obtenga a qué redshifts corresponden las eras de dominio de la radiación, materia y constante cosmológica, y compruebe que la recombinación del hidrógeno se produce durante la era de dominio de la materia. Calcule también la edad del universo en la actualidad y en el instante de recombinación.

\section*{Problema 3}
Teniendo en cuenta que el promedio térmico del producto de la sección eficaz del proceso $e+p\rightarrow H+\gamma$ por la velocidad de los electrones libres es
\begin{equation}
  \langle\sigma v\rangle\simeq 13.6\times10^{-14}\,(T_{\text{eV}})^{-1/2}\,\text{cm}^3\,\text{s}^{-1},
\end{equation}
determine qué fracción de los electrones que quedaron libres después de la recombinación del helio permanecen libres al final de la recombinación del hidrógeno.

\section*{Problema 4}
Teniendo en cuenta que la sección eficaz del scattering Thomson con electrones libres es
\begin{equation}
  \sigma_T\simeq 6.65\times10^{-25}\,\text{cm}^2,
\end{equation}
obtenga la temperatura a partir de la cual la luz se propaga libremente sin interactuar con la materia.

\section*{Problema 5}
Explique por qué la recombinación del hidrógeno no es exactamente un proceso de equilibrio, y por qué la hipótesis de equilibrio da lugar a una sobreestimación de la temperatura de recombinación.

\end{document}
